\scanpage{21}
% archive foliation: 20
\[
\xi'''\in\operatorname{Hom}(N,J)/\operatorname{Im}\operatorname{Hom}(M_{v_0},J)
\]
\deleted{\unsure{Cet}} Cet élément localisé en $\eta$ est nul, \deleted{\unsure{et}} et
c'est là tout ce que nous pouvons tirer du
fait que $u_0\circ\varphi_0$ et $u_{0\eta}$ se prolongent. Mais
\uline{nous savons plus} : les prolongements \inserted{$v$ et $u_\eta$} peuvent se faire
de façon que $v_\eta=\deleted{\unsure{\varphi_\eta}}u_\eta\varphi_\eta$. À priori,
\deleted{\unsure{l'on}} une fois acquis le fait de la possibilité
\uline{séparée} des deux \unsure{constructions|obstructions}, i.e. dans $(1)$
de l'existence de deux éléments dans les
coins, la \inserted{$\unsure{\Rightarrow}$} possibilité de les choisir de
façon cohérente se trouve dans
\[
\operatorname{Hom}(M_{v_0},J_\eta)/\operatorname{Im}
(\operatorname{Hom}(M_{v_0},J)+\operatorname{Hom}(M_{u_0},J_\eta))
\]
\[
\subset\operatorname{Hom}(N,J_\eta)/\operatorname{Im}\operatorname{Hom}(M_{v_0},J).
\]
Soit $\eta$ cet élément. Notons que $\xi'''$ définit
également un élément de cet espace,
puisque
\[
\operatorname{Hom}(N,J_{\unsure{\eta}})/\operatorname{Im}\operatorname{Hom}(M_{v_0},J)
\subset\operatorname{Hom}(N,J_\eta)/\operatorname{Im}\ldots
\]
Nous \unsure{admettrons|démontrons} que l'on a $\eta=\xi'''$. Comme
par hypothèse $\unsure{\eta|\eta'}=0$, on a $\xi'''=0$, OK.
